\chapter{Capturing and Storing Tick Data}\label{ch:pl:capturing-and-storing-tick-data}

The consolidated US options feed is planned for hundreds of billions of messages on a busy day (the dated box of
\cref{sec:pl:capturing-and-storing-tick-data:economics} gives the figure). A firm that keeps every one of those messages for
five years is running a storage business on the side, and the first decision it makes is not which database to buy but what
exactly to keep: the packets as they arrived, with their duplicates, their gaps and their arrival times, or the events they
meant, cleaned, arbitrated and laid out for reading. This chapter captures ten busy minutes of Book~10's simulated exchange
on both of its feed lines, turns them into the firm's event records, measures what each layout costs in bytes, and scales the
result, with cited message counts and prices, to the question of the title of its weekend problem: a petabyte a year?

\section{What tick data is, and where it is captured}

\begin{definition}[Tick data]\label{def:pl:capturing-and-storing-tick-data:tick}
\emph{Tick data}\index{tick data} is the record of every event a venue publishes about its instruments --- order additions,
modifications, cancellations and executions, trades, quotes, status and trading-action messages --- with the timestamps the
venue put on them and the time the firm received them.
\end{definition}

\omterm{def:pl:capturing-and-storing-tick-data:tick}{Tick data} is the finest record of a market a firm can keep, and everything coarser is derived from it: bars (One Quant Book~7,
chapter~2), order-book features (chapter~8 of that book), the replays of its chapter~18. What is less obvious is that ``the''
\omterm{def:pl:capturing-and-storing-tick-data:tick}{tick data} of a day depends on where it was captured. \Cref{fig:pl:capturing-and-storing-tick-data:points} shows the three
usual capture points on the path of Book~13's market-data handling.

\begin{omfigure}
\begin{tikzpicture}[font=\scriptsize,>=Stealth,
  box/.style={draw,rounded corners=2pt,fill=omDef!10,minimum height=0.7cm,inner sep=3pt,align=center},
  cpt/.style={draw,rounded corners=2pt,fill=omThm!12,minimum height=0.6cm,inner sep=3pt,align=center}]
\node[box] (a) at (0,0.6) {line A};
\node[box] (b) at (0,-0.6) {line B};
\node[box] (fh) at (3.2,0) {feed handler\\arbitration, gaps};
\node[box] (bk) at (6.5,0) {book builder,\\ strategy};
\node[box] (st) at (10.0,0) {conflated view\\(screens, risk)};
\draw[->] (a) -- (fh);
\draw[->] (b) -- (fh);
\draw[->] (fh) -- node[above]{events} (bk);
\draw[->] (bk) -- node[above]{updates} (st);
\node[cpt] (c1) at (1.4,-1.9) {1. at the wire:\\packets, both lines};
\node[cpt] (c2) at (4.9,-1.9) {2. after arbitration:\\one event per message};
\node[cpt] (c3) at (8.3,-1.9) {3. after conflation:\\what was acted on};
\draw[->,dashed,omThm] (1.4,0) -- (c1);
\draw[->,dashed,omThm] (4.9,0) -- (c2);
\draw[->,dashed,omThm] (8.3,0) -- (c3);
\end{tikzpicture}

\omcaption{Three capture points on the market-data path. At the wire (a network tap or capture appliance, One Quant Book~14,
chapter~5) the packets of both lines are kept as they arrived; after the feed handler's arbitration each message appears once,
in sequence; after conflation only the updates a consumer actually saw remain.}\label{fig:pl:capturing-and-storing-tick-data:points}
\end{omfigure}

Each point answers different questions. Only the wire capture can say which line lost which packet and when each packet
arrived, which is what a latency or network investigation needs. Only the capture after arbitration is the market as the firm's
book saw it, one message at a time. Only the capture after conflation is what a slow consumer (One Quant Book~13, chapter~12)
actually acted on. Research wants the second; operations and compliance sometimes need the first or the third.

\begin{definition}[Raw capture]\label{def:pl:capturing-and-storing-tick-data:raw}
A \emph{raw capture}\index{raw capture} is the sequence of bytes received at a capture point, kept without interpretation,
each unit (a packet, a frame) with the time the firm received it.
\end{definition}

The chapter's capture is at the first point, in the recorded-file format of Book~10's simulator: for every packet received on a
line, its receive time (8 bytes), its length (4 bytes) and the MoldUDP64 packet itself. Ten minutes of the simulator's busy open,
two instruments, with its seeded impairments --- independent losses on each line, bursts of loss, duplicated packets, jitter and
a five-second outage on line B --- give 15.8~MB on line A and 15.6~MB on line B.

\section{Raw capture against normalised storage}

To turn the capture into the market the book saw, the two lines are merged by sequence number: the first copy of each message
wins, later copies are duplicates, and a gap is held open until the other line has had time to fill it. What neither line fills
is a gap in the record; a live feed handler would ask the venue's retransmission service or rebuild from a snapshot (One Quant
Book~13, chapter~18), and a capture that is written for research records it and moves on.
\Cref{lst:pl:capturing-and-storing-tick-data:arb} is the arbitration of \texttt{firm.tickcap}.

\omcode{code/firm/tickcap/firm_tickcap.py}{53}{87}{Line arbitration for storage: first copy wins, duplicates are counted, and a
gap that neither line fills within the timeout is recorded and skipped.}\label{lst:pl:capturing-and-storing-tick-data:arb}

On the ten-minute session the two lines carry 503\,024 packets. Arbitration delivers 258\,535 messages, discards 255\,899
duplicates (nearly every packet arrives twice, once per line) and records four gaps, eight messages in all, lost on both lines at
once. Each delivered message is then normalised into the firm's event record: the receive time followed by the 48-byte event of
Book~13's feed handler (kind, side, instrument, sequence number, exchange time, order references, price, quantity), 56 bytes in
all. The store and the trading path therefore share one layout, and a test checks that the capture's events are, message for
message, the feed handler's.

\begin{proposition}[Raw is sufficient, normalised is not]\label{prop:pl:capturing-and-storing-tick-data:sufficient}
The normalised records of a session are a function of its \omterm{def:pl:capturing-and-storing-tick-data:raw}{raw capture} (and of the normaliser's code). The \omterm{def:pl:capturing-and-storing-tick-data:raw}{raw capture} is not a
function of the normalised records.
\end{proposition}

\begin{proof}
The first part is the construction: arbitration and normalisation are deterministic functions of the captured packets and their
receive times. For the second, two captures that differ only in what normalisation discards --- which line delivered a message
first, the duplicates, the packet boundaries, the receive time of a later copy, the bytes of a message type the normaliser does
not know --- give the same records.
\end{proof}

The proposition is the whole argument for keeping \omterm{def:pl:capturing-and-storing-tick-data:raw}{raw captures}. A normaliser has bugs, and a field that nobody needed when it was
written (a flag, a new message type after a venue's upgrade) is needed a year later; with the \omterm{def:pl:capturing-and-storing-tick-data:raw}{raw capture} both are fixed by
running the corrected normaliser again, without it they are lost. The price is the bytes, which the rest of this chapter
counts.

\begin{example}[A field nobody kept]\label{ex:pl:capturing-and-storing-tick-data:field}
The simulator's executed-with-price messages (type C) carry a printable flag that says whether the execution counts toward the
day's volume, and its order additions carry the instrument's symbol and a tracking number. The normalised record has no field for
any of them. A researcher who later wants volume that excludes non-printable executions can recompute
it only from the \omterm{def:pl:capturing-and-storing-tick-data:raw}{raw capture}; from the normalised store, the information is gone for every day already written.
\end{example}

\section{Record layouts and encodings}

\Cref{fig:pl:capturing-and-storing-tick-data:layouts} counts bytes per message. The \omterm{def:pl:capturing-and-storing-tick-data:raw}{raw capture} costs 61.0 bytes a message on one
line: the simulator sends one message per packet, so each message carries a 12-byte file record header, a 20-byte MoldUDP64
header and a 2-byte length besides its own 23 to 44 bytes. Real feeds pack many messages into a packet and pay less framing;
the share is a property of the feed, and it is one reason to measure on the feed one actually captures. The normalised record is
56 bytes, fixed: it is barely smaller, but every record is at a known offset, which chapter~4 exploits.

\begin{definition}[Compression ratio]\label{def:pl:capturing-and-storing-tick-data:ratio}
The \emph{compression ratio}\index{compression ratio} of a layout is the size of the uncompressed data divided by the size of the
compressed data, $\mathrm{CR} = b_{\mathrm{raw}} / b_{\mathrm{comp}}$; bytes per message after compression are
$b_{\mathrm{raw}} / \mathrm{CR}$.
\end{definition}

\begin{omfigure}
\begin{tikzpicture}
\begin{axis}[width=12cm,height=5.4cm,xbar,bar width=9pt,xmin=0,xmax=72,
  ytick={0,1,2,3,4,5},yticklabels from table={figdata/platforms/02-capturing-and-storing-tick-data/layouts.csv}{layout},
  xlabel={bytes per message},tick label style={font=\small},label style={font=\small},y dir=reverse,
  nodes near coords,nodes near coords style={font=\footnotesize,/pgf/number format/fixed,/pgf/number format/precision=1},
  grid=major,grid style={gray!20},table/col sep=comma]
\addplot[fill=omDef!45,draw=omDef] table[x=bytes_per_msg,y=k]{figdata/platforms/02-capturing-and-storing-tick-data/layouts.csv};
\end{axis}
\end{tikzpicture}

\omcaption{Bytes per message of the ten-minute session by layout: the \omterm{def:pl:capturing-and-storing-tick-data:raw}{raw capture} of one line, the normalised 56-byte records,
and each compressed by a generic compressor (zlib at level 6), in a columnar file with zstd (chapter~3), and as columns with the
time, sequence and price columns delta-encoded before zlib. Data: \texttt{pl\_tickcap.study}.}\label{fig:pl:capturing-and-storing-tick-data:layouts}
\end{omfigure}

A generic compressor on the \omterm{def:pl:capturing-and-storing-tick-data:raw}{raw capture} gets a ratio of 3.32; on the normalised records, 3.67. The columnar layouts do better,
and the best of the chapter is the simplest columnar trick: write each field as its own column, and write the columns whose
values increase slowly (receive time, sequence number, exchange time) or move by small steps (price) as differences from the
previous value, before a generic compressor. The deltas of the sequence number are all 1, those of the receive time are small
integers, and the ratio rises to 4.70, 11.9 bytes a message. Chapter~3 explains why columns compress better than rows; this
chapter only counts.

\omcode{code/firm/tickcap/firm_tickcap.py}{121}{131}{Delta encoding by column: each field written as a column, the slowly
changing ones as differences from the previous row.}\label{lst:pl:capturing-and-storing-tick-data:delta}

\begin{remark}[The message mix]\label{rem:pl:capturing-and-storing-tick-data:mix}
Of the session's 258\,535 messages, 49.3\% are order additions and 47.3\% deletions; executions are 3.4\%. That is the shape of
a quoted market, and it is why a trades-only store is a hundred times smaller than a full-depth one and answers none of the
questions of One Quant Book~7's chapters~8 and~18.
\end{remark}

\section{Partitioning and compression}

\begin{definition}[Partitioning, partition key]\label{def:pl:capturing-and-storing-tick-data:partition}
\emph{Partitioning}\index{partitioning} splits a dataset into separately stored parts by the value of one or more fields, its
\emph{partition key}\index{partition key} (the trading date, the instrument), so that a query that names key values reads only
the parts that carry them.
\end{definition}

\begin{proposition}[Pruning]\label{prop:pl:capturing-and-storing-tick-data:pruning}
A query whose filter fixes the \omterm{def:pl:capturing-and-storing-tick-data:partition}{partition key} to a set $S$ of values reads $\sum_{v\in S} B_v$ bytes, where $B_v$ is the size of
the partition of value $v$, instead of the whole dataset; with per-file overhead $o$ it also pays $|S|\,o$.
\end{proposition}

\begin{proof}
Partitions outside $S$ cannot contain a matching row, so they need not be opened; each opened partition costs its size and its
overhead.
\end{proof}

The chapter's session is written under \texttt{date=2026-09-28/locate=N/}. Its two instruments are very unequal: instrument~1
takes 811\,160 bytes of the 14\,477\,960, 5.6\% of the day, so a query about it reads 5.6\% of what the unpartitioned file would
cost. The overhead term is the limit. Partition a US options day by option series and a query for one underlying opens
thousands of small files, each costing a directory lookup and an open; partition by underlying and each file is large and the
query opens few. The rule of thumb that follows --- choose the \omterm{def:pl:capturing-and-storing-tick-data:partition}{partition key} from the filters queries actually use, and keep
partitions large enough that their overhead is negligible --- is the subject of chapter~4's tick store.

Compression has its own trade-off, between ratio and speed (\cref{fig:pl:capturing-and-storing-tick-data:speed}). On this laptop
zlib at level 9 gains about 2\% of ratio over level 6 at a twelfth of its speed, and zstd at level 9 inside a columnar file is
both faster and better than zlib at level 6 on the same records. Capture writes once and in real time, so it wants speed; the historical store writes once
at night and is read thousands of times, so it wants ratio.

\begin{omfigure}
\begin{tikzpicture}
\begin{axis}[width=11cm,height=6.8cm,xmode=log,xlabel={compression throughput (MB/s, log scale)},ylabel={compression ratio},
  tick label style={font=\small},label style={font=\small},grid=major,grid style={gray!20},xmin=1.5,xmax=1000,ymin=2.7,ymax=5.0,
  nodes near coords,point meta=explicit symbolic,visualization depends on={value \thisrow{anchor}\as\myanchor},
  nodes near coords style={font=\scriptsize,anchor=\myanchor},table/col sep=comma]
\addplot[only marks,mark=*,omDef] table[x=mb_per_s,y=ratio,meta=method]{figdata/platforms/02-capturing-and-storing-tick-data/measured_compress.csv};
\end{axis}
\end{tikzpicture}

\omcaption{\omterm{def:pl:capturing-and-storing-tick-data:ratio}{Compression ratio} against throughput for the session's \omterm{def:pl:capturing-and-storing-tick-data:raw}{raw capture} (``raw'') and normalised records, by method and
level, best of five runs on one core. Measured on a laptop (Intel Core Ultra 7 155H) under WSL2, machine otherwise idle. Data:
\texttt{bench\_compress.py}.}\label{fig:pl:capturing-and-storing-tick-data:speed}
\end{omfigure}

\begin{method}[Choosing a layout for a capture]\label{met:pl:capturing-and-storing-tick-data:choose}
Keep the \omterm{def:pl:capturing-and-storing-tick-data:raw}{raw capture} of the lines the firm trades on, compressed fast, for as long as investigations need it. Normalise into
fixed-width records in the order of arrival for the intraday store. At the end of the day, rewrite into columns, partitioned by
date and by the key queries filter on, compressed for ratio. Measure bytes per message and the compressor's speed on the
firm's own feeds, not on a benchmark's.
\end{method}

\section{Retention, tiers and the economics of petabytes}\label{sec:pl:capturing-and-storing-tick-data:economics}

\begin{definition}[Storage tier, retention policy]\label{def:pl:capturing-and-storing-tick-data:tier}
A \emph{storage tier}\index{storage tier} is a class of storage with its own price, access time and access charge (hot storage
read in milliseconds, infrequent-access storage, archive storage restored in hours). A \emph{retention policy}\index{retention policy}
states how long each dataset is kept, in which tier at each age, and when it is deleted.
\end{definition}

\begin{dated}{2026-09}{A full feed's message count and storage prices}\label{dat:pl:capturing-and-storing-tick-data:numbers}
The operator of the US consolidated options feed projects, in a notice to subscribers of 15 September 2025 (consulted September
2026), capacity for 311 billion messages a day from July 2026, with peaks of 13.6 million messages per 100 milliseconds, for one
of the two redundant streams. A major cloud provider's list prices for object storage in its US East region (consulted September
2026) are USD 0.023 per GB-month for standard storage, 0.0125 for infrequent access and 0.00099 for its deepest archive tier.
\end{dated}

\begin{dated}{2026-09}{How long order records are kept}\label{dat:pl:capturing-and-storing-tick-data:retention}
In the European Union (consulted September 2026), investment firms engaged in algorithmic trading keep the records of their
orders for five years from the submission of each order (Commission Delegated Regulation (EU) 2017/589, article 28). Market data
is not an order record, but the data a firm needs to explain an order is kept at least as long.
\end{dated}

The storage model of \texttt{firm.tickcap} multiplies it out: bytes a day are messages times bytes per message times the number
of copies; a year is 252 trading days; the steady-state cost of keeping five years is the last month (21 trading days) in hot
storage, the rest of the first year in infrequent access and the four older years in the archive.
\Cref{tab:pl:capturing-and-storing-tick-data:economics} applies it to the options feed's planned capacity with the session's
bytes per message.

\begin{table}[ht]
\centering\small
\begin{tabular}{lrrrr}
\toprule
layout & TB a day & PB a year & PB in 5 years & USD a year, tiered \\
\midrule
\omterm{def:pl:capturing-and-storing-tick-data:raw}{raw capture}, both lines & 37.95 & 9.56 & 47.8 & 1\,989\,000 \\
\omterm{def:pl:capturing-and-storing-tick-data:raw}{raw capture}, one line & 18.97 & 4.78 & 23.9 & 995\,000 \\
normalised & 17.42 & 4.39 & 21.9 & 913\,000 \\
delta columns, zlib 6 & 3.70 & 0.93 & 4.67 & 194\,000 \\
\bottomrule
\end{tabular}
\caption{Full capture of the options feed at its planned daily capacity (dated box), with the simulated session's bytes per
message: storage a day, a year and after five years, and the yearly cost of keeping five years in three tiers at the dated
prices. Data: \texttt{pl\_tickcap.economics}.}\label{tab:pl:capturing-and-storing-tick-data:economics}
\end{table}

\begin{omfigure}
\begin{tikzpicture}
\begin{axis}[width=11cm,height=5.6cm,xlabel={years kept},ylabel={petabytes stored},tick label style={font=\small},
  label style={font=\small},grid=major,grid style={gray!20},xmin=0,xmax=5,ymin=0,ymax=50,xtick={0,1,2,3,4,5},
  legend columns=2,legend style={font=\footnotesize,at={(0.5,-0.28)},anchor=north,draw=none,/tikz/every even column/.append style={column sep=8pt}},
  table/col sep=comma]
\addplot[thick,omThm] table[x=years,y=s0]{figdata/platforms/02-capturing-and-storing-tick-data/stored.csv};
\addplot[thick,omThm,dashed] table[x=years,y=s1]{figdata/platforms/02-capturing-and-storing-tick-data/stored.csv};
\addplot[thick,omDef] table[x=years,y=s2]{figdata/platforms/02-capturing-and-storing-tick-data/stored.csv};
\addplot[thick,omMeth] table[x=years,y=s3]{figdata/platforms/02-capturing-and-storing-tick-data/stored.csv};
\legend{raw; both lines,raw; one line,normalised,delta columns + zlib}
\end{axis}
\end{tikzpicture}

\omcaption{Petabytes accumulated by full capture of the options feed at its planned capacity, by layout: five years of \omterm{def:pl:capturing-and-storing-tick-data:raw}{raw
capture} of both lines reach 48~PB, five years of compressed columns 4.7~PB. Data: \texttt{pl\_tickcap.economics} (message count
and prices in the dated box).}\label{fig:pl:capturing-and-storing-tick-data:stored}
\end{omfigure}

Two cautions keep the table honest. The message count is a capacity the feed's operator plans for, a ceiling rather than an
average day, so the table is an upper bound. And the bytes per message are the simulator's: real options messages have other
sizes and other redundancy, and a real ratio can be higher or lower. The method transfers; the digits must be re-measured on the
feed itself. What does transfer is the shape: the cheapest petabyte is the one not written twice (both lines), not written in
its heaviest layout, and not kept in the hot tier after the month in which anyone reads it.

\section{Tutorial: ten minutes on two lines}\label{tut:pl:capturing-and-storing-tick-data:tut}

\begin{tutorial}
\textbf{Goal.} Capture a busy simulated session on both lines, arbitrate and normalise it, measure every layout, and scale the
result to a full feed. \textbf{End state:} \cref{fig:pl:capturing-and-storing-tick-data:layouts,fig:pl:capturing-and-storing-tick-data:stored}
and \cref{tab:pl:capturing-and-storing-tick-data:economics}.
\begin{tutsteps}
\item \textbf{Generate.} \texttt{pl\_tickcap.day()} runs Book~10's \texttt{make\_recorded\_day} for ten minutes and two
instruments into \texttt{data/platforms/generated/} (about 43~MB, git-ignored, regenerated by the script).
\item \textbf{Arbitrate and normalise} (\cref{lst:pl:capturing-and-storing-tick-data:arb}): 258\,535 messages, 255\,899
duplicates, four gaps of eight messages.
\item \textbf{Measure} with \texttt{study()}: bytes per message and \omterm{def:pl:capturing-and-storing-tick-data:ratio}{compression ratio} of each layout.
\item \textbf{Partition} with \texttt{firm\_tickcap.partition} and compare the bytes a one-instrument query reads.
\item \textbf{Scale} with \texttt{economics()} and the cited inputs.
\end{tutsteps}
\textbf{What to change next.} Sort the records by instrument before the delta encoding and predict, then measure, the change in
ratio; run \texttt{bench\_compress.py} with zstd at more levels and find the level at which the capture would keep up with the
feed's peak rate.
\end{tutorial}

\section{Build: the tick capture}\label{bld:pl:capturing-and-storing-tick-data:tickcap}

\begin{build}
\bfield{Purpose} The first stage of the firm's data spine: \omterm{def:pl:capturing-and-storing-tick-data:raw}{raw captures} of each line, arbitrated normalised records, a
partitioned layout, and the numbers that size the storage. Chapter~4's tick store reads its records; chapter~7's quality checks
read its gaps and counters.
\bfield{Interface} \texttt{RECORD} (56 bytes: receive time and Book~13's feed event); \texttt{read\_capture(data)};
\texttt{arbitrate(lines, timeout\_ns) -> (messages, gaps, counters)}; \texttt{normalise(messages)};
\texttt{partition(rec, root, date)}, \texttt{read\_partition}; \texttt{delta\_encode}, \texttt{compressed\_size(data, method,
level)}; \texttt{StorageModel(msgs\_per\_day, bytes\_per\_msg, ratio, days\_per\_year, copies)} with \texttt{bytes\_per\_day},
\texttt{bytes\_per\_year}, \texttt{stored\_after}, \texttt{cost\_per\_year}.
\bfield{Rules} \omterm{def:pl:capturing-and-storing-tick-data:raw}{Raw capture} is never modified; arbitration is deterministic (first copy by receive time, then by line); every gap
is recorded with its sequence range; the normalised record is Book~13's event plus the receive time, byte for byte.
\bfield{Acceptance tests} \texttt{code/firm/tickcap/tests/}: the record layout; a capture round trip; arbitration with a loss on
one line, duplicates and a gap on both; the same events as Book~13's feed handler on a recorded session up to the first gap;
partitions written and read back; every compression method smaller than the input and the delta encoding invertible; the
storage model by hand.
\bfield{Stretch} Retransmission and snapshot recovery during capture (reuse \texttt{firm.feedhandler}); a capture that writes
compressed blocks with an index; the model with a message count that grows each year.
\end{build}

\begin{omsources}
\item Options Price Reporting Authority, \emph{Revised OPRA Capacity Projections}, notice to multicast data subscribers,
15 September 2025.
\item Y.~Collet and M.~Kucherawy (ed.), \emph{Zstandard Compression and the application/zstd Media Type}, RFC~8878, 2021;
P.~Deutsch and J.-L.~Gailly, \emph{ZLIB Compressed Data Format Specification}, RFC~1950, 1996.
\item Commission Delegated Regulation (EU) 2017/589 (RTS~6), article~28.
\end{omsources}

\section{Exercises}

\begin{exercise}[$\star$]\label{exo:pl:capturing-and-storing-tick-data:1}
Line A's capture is 15\,773\,858 bytes for 258\,535 messages, one message per packet. How many bytes per message is that, and
what share of it is framing (the 12-byte file record header, the 20-byte packet header and the 2-byte length)?
\end{exercise}

\begin{exercise}[$\star$]\label{exo:pl:capturing-and-storing-tick-data:2}
The normalised records compress with zlib at level 6 to 15.24 bytes a message. What is the \omterm{def:pl:capturing-and-storing-tick-data:ratio}{compression ratio}?
\end{exercise}

\begin{exercise}[$\star$]\label{exo:pl:capturing-and-storing-tick-data:3}
Instrument 1's partition holds 811\,160 bytes of the day's 14\,477\,960. How many records is that, and what share of the day
does a query for instrument~1 read with and without \omterm{def:pl:capturing-and-storing-tick-data:partition}{partitioning} by instrument?
\end{exercise}

\begin{exercise}[$\star\star$]\label{exo:pl:capturing-and-storing-tick-data:4}
Keeping both lines' \omterm{def:pl:capturing-and-storing-tick-data:raw}{raw captures} doubles their cost. Give one question only the two \omterm{def:pl:capturing-and-storing-tick-data:raw}{raw captures} can answer, and one the
arbitrated records answer equally well.
\end{exercise}

\begin{exercise}[$\star\star$]\label{exo:pl:capturing-and-storing-tick-data:5}
At 311 billion messages a day and 11.9 bytes a message, how many terabytes a day and petabytes a year (252 days) does the
compressed columnar layout take?
\end{exercise}

\begin{exercise}[$\star\star$]\label{exo:pl:capturing-and-storing-tick-data:6}
With the dated prices, what does keeping five years of the compressed layout cost per year if everything stays in standard
storage, compared with the tiered policy of \cref{tab:pl:capturing-and-storing-tick-data:economics}?
\end{exercise}

\begin{exercise}[$\star\star\star$]\label{exo:pl:capturing-and-storing-tick-data:7}
\emph{Coding.} Sort the session's records by instrument and then sequence number before delta encoding. Predict the change in
\omterm{def:pl:capturing-and-storing-tick-data:ratio}{compression ratio}, then measure it with \texttt{compressed\_size}, and explain the result.
\end{exercise}

\begin{exercise}[$\star\star\star$]\label{exo:pl:capturing-and-storing-tick-data:8}
\emph{Find the flaw.} ``Our capture compresses every packet with zlib at level 9 before writing it, because that gives the best
ratio on our benchmark.''
\end{exercise}

\section{Problem: A Petabyte a Year?}

\begin{problem}[{Weekend problem --- sizing the capture of a full feed}]\label{pb:pl:capturing-and-storing-tick-data:1}
The chapter's ten-minute session, its layouts, and the options feed's planned capacity and the storage prices of the dated
boxes.

\medskip\noindent\textbf{Part I --- The capture.}
\begin{enumerate}
\item How many packets do the two lines carry, and how many messages does arbitration deliver?
\item How many duplicates are discarded, and why is the number close to half the packets?
\item How many messages were lost on both lines, and what would a live feed handler have done about them?
\item Why can the normalised records not be turned back into the \omterm{def:pl:capturing-and-storing-tick-data:raw}{raw capture}?
\item Name one field of the simulator's messages that the normalised record drops.
\end{enumerate}

\medskip\noindent\textbf{Part II --- Layouts.}
\begin{enumerate}[resume]
\item Give bytes per message for the \omterm{def:pl:capturing-and-storing-tick-data:raw}{raw capture}, the normalised records, and the delta-encoded columns with zlib.
\item What are the \omterm{def:pl:capturing-and-storing-tick-data:ratio}{compression ratios} of the \omterm{def:pl:capturing-and-storing-tick-data:raw}{raw capture} with zlib and of the delta-encoded columns?
\item Why does delta encoding help the sequence number most?
\item Which layout would you keep for how long?
\item Why must the bytes per message be re-measured on the real feed?
\end{enumerate}

\medskip\noindent\textbf{Part III --- Scale.}
\begin{enumerate}[resume]
\item At the planned capacity, how many terabytes a day does \omterm{def:pl:capturing-and-storing-tick-data:raw}{raw capture} of both lines take?
\item How many petabytes does it accumulate in five years?
\item Same questions for the delta-encoded columns.
\item What is the yearly cost of keeping five years of each in the three tiers?
\item Why is the capacity figure an upper bound for a typical day?
\end{enumerate}

\medskip\noindent\textbf{Part IV --- The verdict.}
\begin{enumerate}[resume]
\item State the \emph{named result}: the petabytes a year and the yearly cost of full capture in the heaviest and the lightest
layout.
\item Does the answer to the title's question depend more on the layout or on the tiering?
\item What does \omterm{def:pl:capturing-and-storing-tick-data:partition}{partitioning} by instrument save a query for one instrument on the session, and what does it cost when there
are millions of series?
\item Where in the pipeline should the fast compressor run, and where the slow one?
\item In one sentence: what should a firm keep of its \omterm{def:pl:capturing-and-storing-tick-data:tick}{tick data}, and in which form?
\end{enumerate}
\end{problem}

\section{Interview questions}

\begin{interviewq}[$\star$ \iqroles{developer, researcher}]\label{iq:pl:capturing-and-storing-tick-data:1}
What is the difference between raw and normalised \omterm{def:pl:capturing-and-storing-tick-data:tick}{tick data}, and why would a firm keep both?
\end{interviewq}

\begin{interviewq}[$\star$ \iqroles{developer}]\label{iq:pl:capturing-and-storing-tick-data:2}
Two redundant feed lines are captured. How do you turn them into one stream of messages, and what do you do with a gap?
\end{interviewq}

\begin{interviewq}[$\star\star$ \iqroles{developer}]\label{iq:pl:capturing-and-storing-tick-data:3}
How would you estimate the storage needed to keep five years of a full-depth feed, and what would you measure first?
\end{interviewq}

\begin{interviewq}[$\star\star$ \iqroles{developer}]\label{iq:pl:capturing-and-storing-tick-data:4}
How do you choose a \omterm{def:pl:capturing-and-storing-tick-data:partition}{partition key} for \omterm{def:pl:capturing-and-storing-tick-data:tick}{tick data}, and what goes wrong if you choose one that is too fine?
\end{interviewq}

\begin{interviewq}[$\star\star$ \iqroles{developer, researcher}]\label{iq:pl:capturing-and-storing-tick-data:5}
Why does market data compress well, and which fields compress best?
\end{interviewq}

\begin{interviewq}[$\star\star\star$ \iqroles{developer}]\label{iq:pl:capturing-and-storing-tick-data:6}
Design the tick-data storage of a firm trading US options: what is captured where, in which layout, kept how long, in which
tier, and at what cost order of magnitude?
\end{interviewq}
